
\section{Topological defect structure of TGP bodies}
\label{sec:topological_defect}

\subsection{Motivation: what is a ``body'' in TGP?}

In the $N$-body formulation of TGP, each body is modelled as a
\emph{point source} coupling to the scalar field $\Phi$.  This is
analogous to classical electrodynamics, where electrons are treated as
point charges without deriving their internal structure.  Scenario~C
(the most ambitious interpretation) asks whether a \emph{body} in TGP
is actually a \textbf{topological defect} of the $\Phi$ field:
a localised region where $\Phi$ departs from the vacuum value
$\Phi_0$ and approaches $\Phi\to 0$ (the ``absolute nothing'' sector
$\mathcal{S}_0$).

\subsection{Static spherically symmetric ODE}

We seek a static, spherically symmetric profile
$\Phi(r) = \Phi_0\,g(r)$, with $g(r)\to 1$ as $r\to\infty$ (vacuum)
and $g(0)=g_0<1$ (defect core).  With the kinetic function
$f(g)=1+2\alpha\ln g$ (the TGP nonlinear kinetic term, $\alpha=2$)
and the potential
$V(\Phi_0 g)/\Phi_0^2 = \tfrac{\beta}{3}g^3 - \tfrac{\gamma}{4}g^4$,
the Euler--Lagrange equation gives

\begin{equation}
  f(g)\!\left[g'' + \frac{2}{r}g'\right]
  + \frac{\alpha}{g}(g')^2
  = g^2(1-g),
  \qquad \beta=\gamma=1.
\label{eq:defect_ode}
\end{equation}

The right-hand side $g^2(1-g)>0$ for $g\in(0,1)$, so $g''(0)>0$:
solutions rise from $g_0$ toward $g=1$.  The \textbf{ghost boundary}
\begin{equation}
  g^* = \exp\!\bigl(-\tfrac{1}{2\alpha}\bigr)
      = e^{-1/4} \approx 0.7788
\end{equation}
marks the minimum allowed core depth: for $g_0<g^*$, the kinetic term
$f(g)<0$, rendering the theory unphysical (negative kinetic energy).
The physical domain is $g_0\in(g^*,1)$.

\subsection{Linearised far-field: oscillatory, not Yukawa}

Setting $g=1-\delta$ with $|\delta|\ll 1$ and linearising
Eq.~\eqref{eq:defect_ode}:
\begin{equation}
  \delta'' + \frac{2}{r}\delta' = -\delta
  \quad\Longrightarrow\quad
  (\nabla^2+1)\,\delta = 0.
\label{eq:linearised}
\end{equation}
The general solution regular at the origin is
\begin{equation}
  \delta(r) = A\,\frac{\sin r}{r} + B\,\frac{\cos r}{r}.
\label{eq:oscillatory_tail}
\end{equation}
This is \emph{oscillatory} ($\sin r/r$), \textbf{not} a Yukawa tail
($e^{-r}/r$).  The sign of the mass-squared term in the linearised
field equation is $+1$ (from $V''(\Phi_0)=-\beta<0$ evaluated at the
vacuum), leading to $(\nabla^2+m^2)\delta=0$ with $m^2=\beta>0$.
A Yukawa tail would require $(\nabla^2-m^2)\delta=0$ with $m^2>0$,
which arises only at a \emph{minimum} of $V$.

\textbf{Key result (ex11):} Numerical integration confirms that all
solutions with $g_0\in(g^*,1)$ exhibit oscillatory far-field tails
with amplitude $A\approx0.05$--$0.11$ and residual from a Yukawa fit
$\gtrsim 0.003$.  The full nonlinear $f(g)\ne 1$ ODE does not change
this conclusion.

\subsection{Implication: Path~B as external-source formulation}

Since the free-field TGP vacuum emits oscillatory rather than Yukawa
radiation, the Yukawa profile $\Phi_i(r)=C_i e^{-m_{sp}r}/r$
\emph{cannot} be derived from the vacuum sector of the theory.
It must be imposed as an external source:
\begin{equation}
  (\nabla^2 - m_{sp}^2)\,\Phi = -4\pi\,C_i\,\delta^{(3)}(r-r_i),
\label{eq:external_source}
\end{equation}
whose solution is the Yukawa propagator.  This is \textbf{Path~B}:
the Yukawa ansatz is a \emph{phenomenological input}, not a derivable
consequence of the TGP field equations.

\subsection{Physical analogies}

\begin{enumerate}
\item \textbf{Classical electrodynamics:} the point-charge model
  $(\nabla^2)\Phi=-4\pi q\,\delta^3(r)$ does not derive the electron
  radius; it is an axiomatic external source.  TGP Path~B is
  structurally identical.

\item \textbf{Chameleon/scalar-tensor gravity:} in many modified
  gravity theories, the scalar field couples to matter through an
  external source term; the vacuum propagator alone does not determine
  the particle profile.

\item \textbf{Fuzzy dark matter (FDM):} if $\Phi$ acquires a
  non-trivial background condensate, its density profile could provide
  dark-matter-like contributions.  However, the FDM mass scale
  $m_{boson}\sim10^{-22}$\,eV requires
  $m_{sp}\approx8\times10^{-51}$\,[Planck], which is \emph{below}
  the radiation-suppression threshold $m_{sp}>3.4\times10^{-41}$
  derived in \cite{Efimov1970} (see also the N-body predictions section).
\end{enumerate}

\subsection{Rotation-curve diagnostic (ex14)}

A systematic comparison of TGP Path~B Yukawa against the NGC~3198
rotation curve shows that the Yukawa modification
$G_{\rm eff}(r)=G\exp(-r/\lambda)$ \emph{decreases} the gravitational
force at large $r$, giving velocities \emph{below} the Newton-only
prediction.  The flat observed rotation curve ($v\approx147$\,km/s
at $r=20$\,kpc) requires additional dark mass; TGP Yukawa alone makes
the discrepancy \emph{worse}, not better.  Specifically:

\begin{center}
\begin{tabular}{lcc}
\hline
Model & $v_{\rm rot}(20\,\mathrm{kpc})$ & $\chi^2/N$ \\
\hline
Observed & $147$\,km/s & --- \\
Newton (baryons) & $71$\,km/s & $198$ \\
CDM (ISO halo)   & $197$\,km/s & $84$ \\
MOND             & $\approx138$\,km/s & $58$ \\
TGP $\lambda=100$\,kpc & $64$\,km/s & $220$ \\
TGP $\lambda=10$\,kpc  & $26$\,km/s & $387$ \\
\hline
\end{tabular}
\end{center}

\noindent
\textbf{Conclusion:} TGP (Path~B) is not a dark-matter replacement.
It must coexist with standard CDM, contributing three-body corrections
to the internal structure of dark-matter halos and producing an
observable $G_{\rm eff}(r)$ deviation at scales $r\sim\lambda$.

\subsection{Summary of constraints on $m_{sp}$}

\begin{center}
\begin{tabular}{lll}
\hline
Constraint & Bound on $m_{sp}$ [Planck] & Source \\
\hline
Radiation suppr.\ (LIGO) & $m_{sp}>3.4\times10^{-41}$ & ex13 \\
Gravity range $>$50\,AU  & $m_{sp}<1.1\times10^{-46}$ & ex13 \\
FDM mass scale           & $m_{sp}\approx8\times10^{-51}$ & ex14 \\
\hline
\end{tabular}
\end{center}

The FDM value falls \emph{below} the radiation-suppression bound,
ruling out TGP as an ultralight FDM candidate unless the radiation
coupling is suppressed by an additional symmetry.

The range $3.4\times10^{-41}<m_{sp}<1.1\times10^{-46}$ is
\emph{empty} (the lower bound exceeds the upper), implying that no
single value of $m_{sp}$ simultaneously suppresses TGP radiation and
allows the force to operate at Solar-System scales.

\textbf{Resolution:} TGP with $m_{sp}\gg 1$ (Planck-scale mass)
gives $G_{\rm eff}\approx0$ everywhere except at sub-nuclear
distances and zero radiation for all sub-Planck frequencies.  In
this limit, TGP is a theory of \emph{quantum spacetime foam} at the
Planck scale, decoupled from classical gravity observations entirely.
Classical gravity is then described by GR alone, and TGP provides
quantum corrections only at Planck energies.
